\documentclass[a4paper,11pt]{ltjsarticle}
\usepackage[haranoaji]{luatexja-preset}
\usepackage{amsmath,amssymb}
\usepackage[top=12mm,bottom=10mm,left=14mm,right=14mm]{geometry}
\usepackage{array}
\usepackage{tikz}
\usetikzlibrary{calc}
\pagestyle{empty}
\setlength{\parindent}{0pt}
\newlength{\qcol}\setlength{\qcol}{0.47\textwidth}
\newlength{\qgap}
\newlength{\anscolw}\setlength{\anscolw}{0.30\textwidth}
\newlength{\ansh}
\newcommand{\Q}[2]{\parbox[t]{\qcol}{\raggedright (#1)\hspace{0.6em}$\displaystyle #2$}}
\newcommand{\QR}[4]{\Q{#1}{#2}\hfill\Q{#3}{#4}\par\vspace{\qgap}}
\newcommand{\anscell}[1]{\rule[-\ansdrop]{0pt}{\ansh}(#1)}
\newlength{\ansdrop}

\begin{document}
{\large\bfseries 計算問題　5}\hfill 名前(\underline{\hspace{32mm}})\par\vspace{3mm}
■（1）〜（16）の計算をしなさい。（17）〜（20）は方程式を解きなさい。\par\vspace{3mm}
\setlength{\qgap}{5.4mm}
\QR{1}{10a\div(-2)}{2}{(-3)\times5b}
\QR{3}{(-8)-(-12)}{4}{\dfrac{1}{2}(6x-12)+\dfrac{1}{3}(12x+6)}
\QR{5}{(-36)\div(-4)}{6}{(-5)\times(-13)}
\QR{7}{(-19)+(-21)}{8}{\dfrac{1}{4}(8x+24)-\dfrac{1}{3}(9x+3)}
\QR{9}{-4+3^2+(-1)}{10}{(24a-8)\div4}
\QR{11}{(2x+5)+(3x-7)}{12}{(-7)+(+5)-(-10)}
\QR{13}{1-(-8)\div\left(-\dfrac{4}{5}\right)}{14}{\dfrac{a+5}{2}\times6}
\QR{15}{6\div3\times(-5)}{16}{(3x-7)-(2x-8)}
\QR{17}{\dfrac{3}{4}x+5=\dfrac{7}{6}x}{18}{-5x-13=3x+11}
\QR{19}{5x+2=-8}{20}{0.6x-0.4=5}
\vspace{1mm}
\Q{21}{-2\{2a-5(b-2a)\}}\par\vspace{3mm}
{\small 参考記録　　多田C　　1分46秒}\par
\vspace{2mm}
\setlength{\ansh}{9.5mm}\setlength{\ansdrop}{6mm}%
\begin{center}\begin{tabular}{|p{\anscolw}|p{\anscolw}|p{\anscolw}|}\hline
\anscell{1} & \anscell{2} & \anscell{3} \\\hline
\anscell{4} & \anscell{5} & \anscell{6} \\\hline
\anscell{7} & \anscell{8} & \anscell{9} \\\hline
\anscell{10} & \anscell{11} & \anscell{12} \\\hline
\anscell{13} & \anscell{14} & \anscell{15} \\\hline
\anscell{16} & \anscell{17} & \anscell{18} \\\hline
\anscell{19} & \anscell{20} & \anscell{21} \\\hline
\end{tabular}\end{center}

\end{document}
